% Evaluation section of the revised OJ-CS manuscript (input by main-ojcs.tex). Every number is a macro from numbers.tex.
\section{Evaluation}\label{sec:eval}
\paragraph*{Setup} Cryptographic costs (E1) are measured with one thread (\texttt{RAYON\_NUM\_THREADS=1}) on two platforms. The server platform is a core of an \CpuZen{} (MesoNET Zen): each of three independent campaigns ran with its own seed on its own exclusive node, with the process pinned to one core, and all three used one binary built from commit \texttt{\GitCommit}, whose hash is recorded with the results. A uniformly slower first run of one campaign was excluded by a rule fixed before its rerun and rerun on another node (supplementary material; both runs are in the artifact). For comparison, the laptop platform is one core of an \CpuMac{}; the laptop was not idle during the measurements, and its load is recorded with the results. Each point has \Warmup{} warm-up runs and at least 30 measured runs (\BenchN{} for presentation and verification). We report medians with 95\,\% confidence intervals from a hierarchical bootstrap, which resamples campaigns and then runs. Campaign medians differ by at most \CampSpread{} for every operation. A script generates every number, table and figure from the result files.

\subsection{Attack suite}
All \NAttacks{} attacks are rejected and all \NChecks{} correctness checks pass; \NAttReach{} attacks reach a verifier (a tool, or anyone checking published data), the issuer, the revocation log or Judge, and \NAttProver{} are refused by the honest prover library, which shows nothing about verifier soundness. The supplementary material classifies them by threat and rejecting party. The checks cover honest presentations at depths 0 to 4 with and without escrow, every member of a tree using its own part of the interval, a strict tool, the grace window, 2-of-3 openings, and tools re-deriving every proof parameter from public labels.

\paragraph*{Malicious provers and delegators} \NAttMPReach{} of the \NAttMP{} malicious-prover attacks bypass every client-side check and are rejected by the tool or the issuer, which rebuild the statement from their own inputs. Tools reject proofs for a weaker statement (amount or time bound 0), for another revealed class, or with a quota statement detached from the signed interval. For \NAttMDel{} malicious delegations, one per attenuation relation (child budget, interval bounds, expiry, scope bit or depth beyond the parent's, delegation counter beyond its quota), the prover's statement for that relation is re-bound to the parent's own value; the issuer, which rebuilds the statement, rejects each one, and a debug issuer that omits only the binding of that relation to the child's commitment accepts it (positive control). These tests show that every relation is bound to the child's committed values; the relation itself rests on the knowledge soundness of the range arguments, which a test cannot exercise.

\paragraph*{Deployment and framing attacks} The issuer rejects a delegation for another epoch, a root interval wider than its bound $Q$, a root beyond its per-principal limit, and a root request whose attribute layout is shifted by an extra scope entry (which would otherwise move every later field and evade $Q$ and tracing). Tools reject a presentation for another epoch, at a time far from their clock, with a nonce they never issued or already consumed, or against an accumulator value older than the grace window. The log rejects unauthorized, forged, replayed or stale revocations, a descendant revoking its ancestor, and identifier 0. Judge rejects a tag that the issuer binds to a principal without the principal's signature, and, with the registry taken from the log, a key that the issuer substitutes for the principal's. Finally, \NBitFlips{} random single-bit flips over the whole presentation are all rejected.

\subsection{Costs (E1)}
Table~\ref{tab:costs} gives each operation for the default layout (\NMsgs{} messages).
\begin{itemize}
\item \emph{Presentation}: \PresMs{}~ms to produce and \VerMs{}~ms to verify on the server core (\VerBytesMs{}~ms including deserialization and subgroup checks), and \PresMacMs{} and \VerMacMs{}~ms on the laptop core. A presentation is \PresB{} bytes; escrow adds \EscExtraB{} bytes.
\item \emph{Delegation}: \DelParentS{}~s for the parent and \DelIssuerMs{}~ms for the issuer, once per sub-agent.
\item \emph{Public data}: a tool needs \PubParamsKB{}~kB of public data. Everything else is derived from labels.
\end{itemize}

\begin{table}[t]\centering\caption{Costs, default layout, one core: median [95\,\% CI] on the server core, median on the laptop core (ms); bytes.}\label{tab:costs}
\footnotesize\setlength{\tabcolsep}{2pt}
\begin{tabular}{@{}lrrr@{}}\toprule Operation & Server & Laptop & Bytes\\\midrule
\inputtab tab_costs.tex \end{tabular}\end{table}

\paragraph*{Where the time goes} Adding the statement groups one by one (table in the supplementary material): The BBS+ proof alone costs \AblBbsPres{}~ms. A pseudonym with a counter checked against the signed interval adds \AblRevAdd{}~ms if the counter is revealed, as in $k$-TAA, and \AblDyAdd{}~ms if it is hidden: hiding the counter costs \AblHideCost{}~ms, spent in the two range proofs of the interval. The cap and expiry ranges add \AblRangeAdd{}~ms, and the five non-membership proofs for chained revocation \AblNmAdd{}~ms. The four range proofs dominate; each uses two 64-bit Bulletproofs++ arguments in the library we use, so aggregating them is the obvious optimization.

\paragraph*{Scaling} Presentation cost does not depend on the size of the mandate, its depth or the revocation list: \MsgPresMin{}--\MsgPresMax{}~ms from \MsgLo{} to \MsgHi{} signed messages, \DepthPresMin{}--\DepthPresMax{}~ms for depths 0 to 4, and \RevPresMin{}--\RevPresMax{}~ms up to \RevMax{} revoked entries (figure and table in the supplementary material). Delegation grows linearly with the number of classes (\DelMsgLoMs{}~ms at \MsgLo{} messages, \DelMsgHiS{}~s at \MsgHi{}); after 100 revocations, each holder refreshes its five witnesses in \RevRefreshMs{}~ms.

\paragraph*{Throughput} Deserializing and verifying a presentation takes \ThrVerMs{}~ms, so one core verifies about \ThrPerCore{} presentations per second. Independent single-threaded verifiers scale almost linearly: \ThrMpRate{} per second on \ThrMpCores{} cores, a parallel efficiency of \ThrMpEff{}. A presentation with one flipped bit is rejected after \ThrBadMs{}~ms, and random bytes of the right length are rejected by the parser in about \ThrGarbageUs{}~$\mu$s.

\subsection{End-to-end on MCP (E2)}
\paragraph*{Setup} One agent client calls three MCP servers (Python SDK \EMcpSdk{}, Python \EPy{}), with cryptography on one laptop thread. There are four transports: stdio, streamable HTTP on the loopback interface, HTTP through a relay that adds 10~ms in each direction, and stateless HTTP, with a new connection and MCP session for every call so that no session identifier links calls.
Each mode and transport makes \ECalls{} calls after \EWarm{} warm-up calls. The whole measurement was repeated \ERuns{} times with different seeds with the machine load recorded; we report the median of the run medians. Before each authorized call the agent fetches a single-use nonce and the tool's time from the tool (\EChalHttp{}~ms over HTTP). All modes use the same two-hop delegation (root, agent, sub-agent) and bind each call to the canonical JSON (with currency), the tool identity and the nonce, except where the mechanism cannot.

\paragraph*{Baselines} (1) A JWT chain with one delegation JWT per hop and a per-call proof-of-possession JWT, in the style of Agentic JWT and DPoP~\cite{goswami2025ajwt,rfc9449}. (2) A Biscuit token with two attenuation blocks and a per-call block~\cite{prakash2026aip}. (3) A macaroon with two delegation caveats and a per-call caveat~\cite{birgisson2014macaroons}; tools must share the root key. (4) Batch one-time tokens: the issuer checks the delegation chain in the clear and issues single-use signed tokens with a proof-of-possession key, as with SD-JWT batch issuance~\cite{rfc9901}. (5) RSA blind-signature tokens (RFC~9474) that commit to an ephemeral key which signs the call, as in Privacy Pass with a proof-of-possession binding~\cite{rfc9474,rfc9578}. We added a status list to the JWT chain and to batch tokens so that they can revoke. ARC and anonymous counting tokens are not implemented.

\begin{table*}[t]\centering\caption{End-to-end MCP call (E2) on a laptop that was not idle (load recorded), medians over 3 runs: agent preparation, tool verification and total per call in ms (streamable HTTP, +10~ms per direction, stateless HTTP); authorization bytes; delegation (two hops) or issuance cost per token; whether a modified call is rejected, revoking the agent stops the sub-agent, colluding tools can link calls; what the issuer learns.}\label{tab:e2e}
\footnotesize\setlength{\tabcolsep}{2.5pt}
\begin{tabular}{@{}lrrrrrrrcccl@{}}\toprule Mode & Prep. & Verify & Total & +10 ms & Stateless & Bytes & Deleg./issuance & Binds & Revokes & Tools link & Issuer learns\\\midrule
\inputtab tab_e2e.tex \multicolumn{12}{@{}l}{\scriptsize $^\dagger$ except by a 2-of-3 escrow opening. Blind signing in pure Python.}\end{tabular}\end{table*}

\paragraph*{Results} Table~\ref{tab:e2e} shows the price of unlinkability.
\emph{Cost}: a local HTTP call costs \EtoeNoneHttp{}~ms without authorization, and every baseline stays at or below \EBaseMaxHttp{}~ms. \sys{} takes \EMMPrep{}~ms on the agent and \EMMVer{}~ms on the tool, for \EMMTotalHttp{}~ms per call (\EMMTotalLo{}--\EMMTotalHi{} across runs) and \EMMTotalDelay{}~ms with the added delay. Its two-hop delegation costs \EMMDelegS{}~s once. On the wire a presentation exceeds its \PresB{} raw bytes (base64 in the JSON \texttt{\_meta} field). For scale, a $k$-TAA-style credential with a revealed counter costs \AblRevPres{}~ms to present (ablation); blind RSA issuance is pure Python. \emph{Stateless sessions}: a fresh MCP session per call, which removes the session identifier as a linking handle, adds about \EStatelessAdd{}~ms per call in every challenge-based mode. \emph{Controls}: every mode except ``none'' rejects replays to the same and to another tool, also with a fresh nonce. Without proof of possession, blind tokens would accept a modified call or currency. Macaroons and blind tokens still accept the sub-agent after its parent is revoked. \emph{What the issuer learns}: with batch tokens, the issuer sees every delegation chain and token identifier. With blind tokens, it sees the chain and how many tokens each root requests. With \sys{}, it sees only when delegations happen.
Only \sys{} also offers a tree-wide quota and, with escrow, accountability.

\subsection{Collusion study (E3)}\label{sec:link}
\paragraph*{What we simulate} E3 measures what channels other than the token still reveal. A simulator generates 1{,}000 and 10{,}000 users calling \LinkTools{} tools with Zipf popularity over \LinkDays{} days, which gives \LinkCalls{} calls at 1{,}000 users.
Each user has favorite tools, interests, a home city, a budget and a daily routine. A task is a burst of 2 to 6 calls; the gaps between calls are resampled from \GapN{} gaps of the \sys{} runs of E4 (median \GapMedS{}~s, 99th percentile \GapNinetyNineS{}~s), and timestamps have millisecond resolution. A share of tasks (\IssPDeleg{}, as in E4) starts by delegating to a sub-agent. Booking and payment calls carry personal fields (city, party size, preferences), search calls carry topics, and a payment repeats to the cent a price returned earlier in the task.
All tools pool their logs.

\paragraph*{Attackers} (a) Identifier matching. (b) Timing: links calls closer than a window, and also runs a session tracker. (c) Sessions, then 30-day clustering on call content. (d) As (c), plus tool sets and exact amounts. The \emph{oracle} assigns each session to the user with the most likely true profile (strong, not an upper bound). In the \emph{named} setting, one tool per task (payment or booking) legitimately learns the user's name, and the coalition attaches the rest of the session, and then the clusters, to that name.

\paragraph*{Metrics} \emph{Profile recovery} is the share of users whose largest cluster holds at least half of their calls with at least 80\,\% purity. \emph{Naming} is the share of calls attributed to the right named person. We run \LinkSeeds{} seeds per cell. All numbers come from a synthetic population: they show which channels link and how strongly, not the rates of a real deployment. ``Minimized'' runs strip personal fields from arguments.

\begin{table}[t]\centering\caption{Collusion of 50 tools (E3), \%, mean of 5 seeds (ranges in the supplementary material): profile recovery by attacker (d) and by the profile oracle at 1{,}000 users, (d) at 10{,}000 users, and naming of calls when one tool per task learns the user's identity, at 1{,}000 and 10{,}000 users.}\label{tab:link}
\footnotesize\setlength{\tabcolsep}{2.5pt}
\begin{tabular}{@{}lrrrrr@{}}\toprule & \multicolumn{3}{c}{Profile recovery} & \multicolumn{2}{c}{Naming}\\ Token format & (d) & Oracle & (d), $10^4$ & $10^3$ & $10^4$\\\midrule
\inputtab tab_link.tex \end{tabular}\end{table}

Table~\ref{tab:link} gives four findings.
\begin{enumerate}
\item \emph{Identifiers matter.} Per-tool pseudonyms behave like a principal identifier once tools pool their logs: \LinkPertoolReid{} of profiles recovered. A revealed counter is also a strong signal: \LinkPtejReid{} by (d) and \LinkPtejOracle{} by the oracle, against \LinkMMReid{} and \LinkMMUpper{} for \sys{}. From the tools' view, \sys{} and one-time tokens give identical logs, and they get identical numbers.
\item \emph{Content still links.} With full arguments, (d) recovers \LinkMMReid{} of profiles (\LinkMMReidLo{}--\LinkMMReidHi{} over seeds). With minimized arguments it recovers \LinkMMMinReid{}, and at most \LinkMMMinSensMax{} across the sensitivity settings listed in the artifact.
\item \emph{A tool that learns the identity names the session.} Timing links the calls of a task, so the coalition names \NameMM{} of \sys{} calls at 1{,}000 users and low load (\NameMMLo{}--\NameMMHi{} over seeds). The share rises to \NameMMEvery{} if every task touches such a tool. Minimizing arguments only brings it to \NameMMMin{}, and five colluding tools already name \CoalFiveName{}. Naming weakens as tasks overlap: from \NameLoadLo{} at \LoadRateMin{} tasks per minute across the coalition to \NameLoadHi{} at \NameLoadHiRate{} (\NameLoadHiUsers{} users) (figure in the supplementary material). The attacker attributes every call, so the rest are named wrong. An issuer in the coalition, which knows when delegations happen (roots are long-lived), adds only \IssNameInc{} points (\IssNameIncPdOne{} if every task delegates).
\item \emph{Delays.} Random per-call delays of up to 10~s defeat profile clustering (\JitDTen{}) but not naming (\JitNameSixHundred{} with up to 600~s).
\end{enumerate}
Calling the identity-learning tool in its own session, at a random time within the next hour, cuts naming to \NameDecHour{}; but an attacker that joins the payment to the earlier quote of the same amount brings it back to \NameDecHourAdapt{}, and only a call that carries no value echoed from the task keeps it at \NameDecHourFree{} (\NameDecHourFreeTenK{} at 10{,}000 users). It stays at \NameDecHourEvery{} (original attacker) when every task needs such a tool. Unlinkable authorization is thus necessary but not sufficient: calls must also carry little, and identity-learning calls must be decoupled from the task and carry no value it echoes.

\subsection{LLM agents in the loop (E4)}\label{sec:e4}
\paragraph*{Setup} A language model served by vLLM~\EFourVllm{} on one H100 GPU solves \EFourTasks{} tasks by calling four MCP servers over streamable HTTP (SDK \EFourSdk{}). There are five templates: (a) pay the cheapest flight and hotel under caps; (b) the same, with the hotel delegated to a sub-agent; (c) as (b) plus a transfer, with the sub-agent revoked after its first payment (correct outcome: a reported failure); (d) a search-heavy task; (e) a five-city survey delegated with a three-search quota (correct outcome: flight paid, partial failure reported).
Each task runs in four modes: none, a JWT chain with status list and per-tree counters at the tools, \sys{}, and \sys{} with escrow. It runs in \EFourCfgs{} configurations: temperature 0 with and without tool latency (lognormal, median \EFourLatMed{}~ms), a replication of the first with identical settings, and temperature 0.7 with three seeds. This gives \EFourRuns{} runs over \EFourModels{} open models and \EFourCalls{} tool calls. Tools are strict verifiers; they canonicalize the raw JSON-RPC arguments themselves, with the same JCS implementation and MCP SDK as the client. Authorization travels in \texttt{\_meta}, so prompts, tool outputs and seeds are identical across modes. When a model stops without finishing, the harness sends up to two ``please continue'' messages (\EFourNudgeMin{}--\EFourNudgeMax{}\,\% of runs, depending on the model, in every mode). Agent and tools run on single cores of an \EFourCpu{} (presentation \EFourPresMs{}~ms, verification \EFourVerMs{}~ms).

\begin{table}[t]\centering\caption{LLM agents (E4), six configurations pooled, \%: success on honest templates (a, b, d); correct outcome of revocation (c) and quota (e) templates; median slowdown of \sys{} (MM) over none (tool latency on); share of task time in MM authorization (temperature 0).}\label{tab:e4}
\scriptsize\setlength{\tabcolsep}{1.8pt}
\begin{tabular}{@{}lrrrrrrrrr@{}}\toprule & \multicolumn{3}{c}{Honest} & \multicolumn{2}{c}{Revoked (c)} & \multicolumn{2}{c}{Quota (e)} & & \\ Model & none & JWT & MM & JWT & MM & JWT & MM & Slow. & Share\\\midrule
\inputtab tab_e4.tex \end{tabular}\end{table}

\paragraph*{Results} Table~\ref{tab:e4} answers four questions.
\begin{itemize}
\item \emph{Does authorization change what honest agents achieve?} By construction the model sees the same transcript in every mode, so outcomes can differ only through rejections or serving noise. The substantive result is that, with a per-task issuer, no honest call was wrongly rejected (\EFourHonestCalls{} calls with \sys{}), and client and tool canonical strings never differed, over \EFourCalls{} calls and \EFourProbeCases{} adversarial argument shapes. The rest is serving noise: rerunning a configuration changes up to \EFourRepMax{} of \EFourRepPairs{} outcomes (only \EFourRepDetModel{} is deterministic), so we use a pragmatic margin of $\pm$\EFourMarginPP{} points, half that churn, fixed after observing it. Equivalence means that the 90\,\% cluster-bootstrap interval (unit: task) lies within the margin, i.e.\ two one-sided tests at 5\,\%; per-configuration results are in the supplementary material. With a permutation test that respects the clustering of units in tasks, no difference to ``none'' on the honest templates is significant (smallest $p=\EFourMinP{}$; Holm over \EFourNTests{} tests). \sys{} is equivalent to ``none'' for \EFourMMEqTsev{} of 5 models at temperature 0.7 (not gpt-oss-20b) and for all 5 when the four distinct configurations are pooled (largest difference \EFourMMMaxDiff{} points); over all modes, \EFourEqTsev{} and \EFourEqAll{} of \EFourEqPairs{} mode--model pairs are equivalent.
\item \emph{Does it stop what it should?} All \EFourRevAtt{} calls by revoked sub-agents and all \EFourQuotaAtt{} calls beyond a sub-agent's quota were rejected under JWT and \sys{}. Without authorization, all \EFourRevNone{} and \EFourQuotaNone{} were accepted. Under \sys{} the honest wallet refused to present. A dishonest wallet that kept presenting after revocation (old accumulator value, or stale witness) and beyond its quota (reused counter) had all \EFourDwSent{} (scripted) and \EFourDwLlmSent{} (five models) such presentations rejected by the tools, with the uniform error. Whether the agent then reports the failure depends on the model: \EFourSafeMin{} to \EFourSafeMax{}\,\% correct outcomes for (c), \EFourQuotaMin{} to \EFourQuotaMax{}\,\% for (e). The small models often claim success after a rejection, with JWT as with \sys{}.
\item \emph{What does it cost?} With tool latency, \sys{} slows honest tasks by a median factor of \EFourSlowMin{} to \EFourSlowMax{} (up to \EFourSlowNoLatMax{} without tool latency), and JWT by at most \EFourJwtSlowMax{}. Authorization takes \EFourShareMin{} to \EFourShareMax{} of task time (JWT: at most \EFourJwtShareMax{}). Per-call proofs dominate; the sub-agent's delegation takes \EFourDelegS{}~s. The share falls with slower models or tools; a wallet in a separate process adds only \EFourIpcMs{}~ms per call.
\item \emph{What does a shared issuer change?} One issuer serves all tasks of a worker and revokes an unrelated mandate every \EFourShEvery{}~s. A presentation built just before a revocation and verified just after it is then rejected: with LLM agents, \EFourShRejMin{} to \EFourShRejMax{} of honest calls, and success falls to \EFourShSuccZeroMin{}--\EFourShSuccZeroMax{}\,\% (per-task issuer: \EFourShBaseMin{}--\EFourShBaseMax{}\,\%). A window of two accumulator values removes the rejections but leaves a revoked mandate usable for two updates (\EFourWtwoDelayS{}~s here, \EFourWtwoDelaySparseS{}~s at one revocation per minute). A wallet retry with the same counter and a fresh nonce rejects none of \EFourRetryCalls{} honest scripted calls, with revocation effective after \EFourRetryDelayMs{}~ms; a time-bounded window accepting the previous value for 1~s rejects none and bounds the delay to \EFourDeltaDelayS{}~s whatever the rate. Both delays exclude propagation from the log, and both remedies assume batched revocations (supplementary material). With Qwen2.5-72B and gpt-oss-20b, both reject no honest call and restore success to \EFourLlmFixSuccMin{}--\EFourLlmFixSuccMax{}\,\%.
\end{itemize}

